% Section 9. Drafted from docs/PAPER-OUTLINE.md section 9 and docs/THEORY.md "The compensation-
% commutativity theorem as the base case". Every named result is a gated, axiom-free theorem in
% coq/NormalizationDerived.v unless noted.

\section{The compensation-commutativity theorem as the base case}\label{sec:base-case}

The single-registry theorem of \cref{sec:main-single} was proved directly, before the framework
existed. This section derives it from the framework's kernels, used unchanged. It is a derivation, not
a prediction: the framework was built after the theorem, with it in view. What it shows is that the
kernels, as they are, reproduce the exact statement, and that each of its conditions is a named
condition of a layer. The setting is \cref{thm:cc}'s: free delivery, terminating repair, and
$N = \rho^{*}$.

\paragraph{The dictionary.}
\begin{itemize}[nosep]
  \item \emph{CC2 is state descent.} From a start $s_0$, every event descends along $N$ on the states
    reachable from $s_0$ if and only if CC2 holds there (\coq{nd_descends_cc2}); globally this is
    \cref{prop:descent-cc2}.
  \item \emph{CC1 is history descent on free delivery.} The governed semantics from the repaired start,
    $t \mapsto \mathrm{grun}(t, N(s_0))$, is invariant under the free-delivery presentation, whose
    generated equivalence is equality up to permutation (\coq{swap_peq_perm}). By the local
    principle (\cref{thm:p1}) this holds exactly when every adjacent swap commutes after every word
    (\coq{nd_history_seeds}), and given state descent it is CC1 on the states reachable from $s_0$
    (\coq{nd_history_cc1}). History descent alone gives CC1 only at the repaired states the governed
    runs reach; state descent carries it to the raw reachable states, where it is needed because an
    event can fire before repair.
  \item \emph{Termination is the progress side,} supplied by the potential $\Phi$.
  \item \emph{Newman joins them.} A peak between an event and a repair step is joined by state descent,
    and a peak between two events by CC1 (\coq{nd_state_peaks_join}, \coq{nd_history_peaks_join}); the
    peak classification of the canonical-execution kernel (\coq{classified_peak_exact}) then gives
    confluence (\coq{nd_confluent_from}).
  \item \emph{The converse is by the same small runs} as in \cref{sec:main-single}
    (\coq{nd_tiny_cc1}, \coq{nd_tiny_cc2}).
\end{itemize}

\paragraph{Proof idea for the dictionary.}
For \coq{nd_descends_cc2}, one repair does not change the repaired form of an invalid state,
$N(\rho(\sigma)) = N(\sigma)$ (\coq{nd_N_rho}, because runs from an empty buffer are deterministic);
descent at the reachable states $\sigma$ and $\rho(\sigma)$ then gives CC2 at $\sigma$. Conversely, CC2
keeps $\gov(e)$ unchanged along the repair chain from a reachable $x$ to $N(x)$, by induction on the
chain (\coq{nd_cc2_chain}), which is descent at $x$. For \coq{nd_history_seeds}, the future-respecting
hypothesis of \cref{thm:p1} holds trivially, since the governed semantics of $t \cdot e$ is $\gov(e)$
applied to that of $t$, so \coq{local_exact_all} reduces invariance to the adjacent swaps. For
\coq{nd_history_cc1}, state descent makes $N(\sigma)$ a governed run from $N(s_0)$ for every reachable
$\sigma$ (\coq{nd_reach_canon}) and gives $\gov(e)(\sigma) = \gov(e)(N(\sigma))$, so the swap after
that run is CC1 at $\sigma$; the converse applies CC1 at the governed runs, which are reachable
(\coq{nd_reach_grun}). The peaks join directly: an event against a repair step meets at
$(\gov(e)(s), B - e)$ by CC2, and two events meet at $(\gov(e_2)(\gov(e_1)(s)), B - e_1 - e_2)$ by CC1.

\begin{theorem}[\coq{nd_derived_from}, \coq{nd_cc_exact}]\label{thm:derived}
For every start $s_0$, every buffer delivered from $s_0$ has a unique normal form if and only if
free-history descent holds from $N(s_0)$ and every event descends on the states reachable from $s_0$.
The three original statements of the single-registry theorem follow word for word.
\end{theorem}

\paragraph{Proof idea.}
For sufficiency, \coq{nd_history_cc1} turns the two descent conditions into CC1 on the states
reachable from $s_0$, and \coq{nd_confluent_from} gives confluence from $(s_0, B)$:
\coq{classified_peak_exact} is Newman's lemma on the reachable configurations, with termination from
$\Phi$ and a classifier under which every nontrivial peak is a state peak or a history peak
(\coq{nd_cls_complete}). Confluence gives unique normal forms. For necessity, \coq{nd_tiny_cc1} and
\coq{nd_tiny_cc2} run a word that reaches $\sigma$ and then the two small buffers; both ends are normal
forms because $\rho^{*}$ reaches a valid state and the buffer is empty, so unique normal forms make
them equal and give CC1 and CC2 at $\sigma$. The dictionary turns these back into the two descent
conditions. The original statements follow by rewriting with the dictionary
(\coq{nd_dictionary_from}); \coq{nd_cc_exact}, for configurations over any reachable state, applies
the start-relative form at that state and uses transitivity of reachability.

Only one proof in the file cites the original theorems: \coq{nd_derived_matches_original}, which shows
that the derived statements are the originals' by the dictionary alone. Every other theorem of the file
remains closed under the global context when the original theorems are turned into admitted axioms in
a scratch copy, so none of them depends on the originals. A small instance where both sides hold shows
the hypotheses are satisfiable (\coq{nd_nonvacuous}).

\paragraph{What it shows and what it does not.}
The base case shows that the framework reproduces the theory's central theorem with its conditions
named by layer: state descent at one event, history descent over sequences, termination as progress,
Newman as the bridge. It does not show reach, because the framework was designed with this theorem in
view; that is what \cref{sec:tests} tests.
